\documentclass[11pt,letterpaper]{article}
\usepackage[T1]{fontenc}
\usepackage[utf8]{inputenc}
\usepackage{lmodern}
\usepackage{amsmath,amssymb,amsthm,mathtools}
\usepackage[margin=1in,headheight=15pt]{geometry}
\usepackage{microtype}
\usepackage{booktabs,longtable,array,tabularx}
\usepackage{xcolor,listings,fancyhdr,enumitem}
\usepackage{hyperref}
\usepackage{bookmark}
\definecolor{accent}{HTML}{224C65}
\definecolor{muted}{HTML}{53616B}
\definecolor{codebg}{HTML}{F4F6F7}
\hypersetup{colorlinks=true,linkcolor=accent,citecolor=accent,urlcolor=accent,
  pdftitle={A Detailed Review of DenestRadicals3 and denest11},
  pdfauthor={OpenAI},pdfsubject={Mathematical correctness, Wolfram Language semantics, and software quality}}
\pagestyle{fancy}
\fancyhf{}
\fancyhead[L]{\small\textsc{Radical denesting: source review}}
\fancyhead[R]{\small September 2026}
\fancyfoot[C]{\thepage}
\renewcommand{\headrulewidth}{0.3pt}
\setlength{\parskip}{4pt}
\setlength{\parindent}{1.2em}
\setlength{\emergencystretch}{2em}
\setcounter{tocdepth}{2}
\newcommand{\code}[1]{\texttt{\detokenize{#1}}}
\newcommand{\src}[1]{\hyperref[app:source]{source lines #1}}
\newcommand{\Q}{\mathbb{Q}}
\newcommand{\C}{\mathbb{C}}
\newcommand{\Z}{\mathbb{Z}}
\DeclareMathOperator{\Arg}{Arg}
\DeclareMathOperator{\Disc}{Disc}
\DeclareMathOperator{\lcm}{lcm}
\DeclareMathOperator{\raddepth}{rdepth}
\newtheorem{proposition}{Proposition}[section]
\newtheorem{lemma}[proposition]{Lemma}
\newtheorem{theorem}[proposition]{Theorem}
\theoremstyle{definition}
\newtheorem{definition}[proposition]{Definition}
\theoremstyle{remark}
\newtheorem{remark}[proposition]{Remark}
\lstdefinelanguage{Wolfram}{
  morekeywords={Block,Module,If,For,While,Which,Return,True,False,Null,
    ClearAll,Unprotect,SetAttributes,BeginPackage,Begin,End,EndPackage,
    Function,Rule,RuleDelayed,Set,SetDelayed,With,Table,Options,OptionsPattern},
  sensitive=true,morecomment=[s]{(*}{*)},morestring=[b]"
}
\lstset{language=Wolfram,basicstyle=\ttfamily\small,
 keywordstyle=\color{accent},commentstyle=\color{muted},stringstyle=\color{muted},
 backgroundcolor=\color{codebg},frame=single,framerule=0.2pt,rulecolor=\color{black!15},
 breaklines=true,breakatwhitespace=false,columns=fullflexible,keepspaces=true,
 showstringspaces=false,tabsize=2,aboveskip=8pt,belowskip=8pt,
 xleftmargin=5pt,xrightmargin=5pt}
\setlist{nosep,leftmargin=1.6em}
\begin{document}
\begin{titlepage}
\vspace*{0.55in}
{\color{accent}\large\textsc{Mathematical software review}\par}
\vspace{0.25in}
{\fontsize{28}{33}\selectfont\bfseries A detailed review of a\par radical-denesting program\par}
\vspace{0.2in}
{\Large\code{DenestRadicals3}, \code{denest11}, and their helpers\par}
\vspace{0.35in}
{\large Algorithm, mathematical justification, defects, and repair strategy\par}
\vspace{0.25in}
{September 5, 2026\par}
\vspace{0.5in}
\noindent\rule{\textwidth}{0.7pt}
\begin{abstract}
The supplied Wolfram Language program attempts to denest algebraic radical
expressions by combining structural preprocessing with a search over
multipliers. Its central experiment computes the minimal polynomial of a
scaled root and intersects that polynomial with a binomial over an algebraic
coefficient field. A lower-degree greatest common divisor may reveal a simpler
representation. This is a mathematically meaningful strategy, and the
program's denominator-rationalization identity is correct on its natural
domain.

The implementation nevertheless has important correctness defects. It can
lose the branch of the original expression, applies unguarded power identities
in its custom factorizer, uses a heuristic zero test as a certificate, and
contains a comparator that ignores its second argument. Its nominal
multiplier cap is not a cap, invalid multipliers can lead to undefined
arithmetic and unchecked indexing, and loading the file changes the protection
of a system symbol. The article distinguishes these defects from heuristic
incompleteness, representation choices, and merely unusual but valid language
idioms. It supplies proofs, exact counterexamples, a source map, a defensive
single-trial reference component, and a reproducibility package.
\end{abstract}
\vfill
\noindent\textbf{Evidence boundary.} This is a complete source inspection
supported by official language documentation and independent algebraic checks.
The supplied program and the proposed Wolfram tests were \emph{not} executed
in a Wolfram kernel. Claims about mathematical rules are proved below;
version-dependent evaluation traces are identified as regression tests rather
than reported observations.
\end{titlepage}
\tableofcontents
\clearpage

\section{Scope, provenance, and overall assessment}
\label{sec:scope}
The object reviewed is the attachment \code{Pasted text.txt}, reproduced
without editorial changes as \code{original.wl} in the accompanying archive.
It contains 585 physical lines and 21,199 bytes. Its SHA-256 digest is
\begin{center}\small
\path{693a2825f1e2b8912e55371b90dc8ed115bc33ee707dd3a4521eb22bcd2ccc93}.
\end{center}
All source-line references in this article refer to that file, not to the
line numbers of the LaTeX source. Appendix~\ref{app:source} gives a numbered
listing. The attachment supplies no test suite, package declaration, supported
kernel version, formal input contract, or proof of completeness. No such
information is inferred from the names of the functions.

\subsection{What the program is trying to do}
A representative objective is to replace
\[
 \sqrt{3+2\sqrt2}\quad\text{by}\quad 1+\sqrt2,
 \qquad
 \sqrt{5+2\sqrt6}\quad\text{by}\quad \sqrt2+\sqrt3.
\]
Both changes preserve the number while reducing the nesting of explicit
radicals. The implementation is more ambitious than a formula for nested
square roots: it accepts rational powers, constructs algebraic multipliers,
uses arbitrary-degree minimal polynomials, and explicitly enumerates roots of
unity. Its architecture is best described as an \emph{experimental,
multiplier-guided algebraic simplifier}, not as a demonstrated complete
denesting algorithm.

The main verdict is not that the algebraic idea is unsound. The main verdict
is that \textbf{a promising search engine is wrapped in transformations and
acceptance criteria that do not reliably preserve the original expression}.
A search may be incomplete and still be useful; an exact simplifier must not
silently return the wrong algebraic number.

\subsection{Three kinds of evidence}
Throughout the article, the following distinctions are maintained.
\begin{description}
\item[Source-established finding.] A property follows directly from the supplied
statements: for example, \code{Unprotect[Print]} is executed at load time and
there is no matching restoration.
\item[Mathematical counterexample or independent check.] An identity, branch
calculation, polynomial, or counting formula can be proved independently of
the original evaluator. These checks were also exercised with Python,
SymPy~1.14.0, and mpmath where appropriate.
\item[Wolfram regression probe.] A call intended for a real Wolfram kernel tests
how automatic evaluation, factoring, pattern matching, and solver choices
interact end to end. Such probes are included but are not presented as
executed tests.
\end{description}
Official Wolfram documentation is used only to establish language behavior,
not to replace the attached algorithm with a different one. In particular,
principal powers, exact algebraic comparison, scoping, factoring options, and
sorting semantics are checked against the documentation
\cite{power,possiblezero,rootreduce,block,sort,factorinteger}.

\section{Prioritized findings}
\label{sec:findings}
The table is a navigation aid. ``Critical'' means an incorrect value can be
accepted on a relevant branch or through an invalid transformation. ``High''
includes uncontrolled failure and potentially severe resource consumption.
The severity of heuristic incompleteness depends on the promised interface;
no completeness promise is present in the attachment.

\small
\begin{longtable}{@{}p{0.075\textwidth}p{0.13\textwidth}p{0.48\textwidth}p{0.235\textwidth}@{}}
\toprule
ID & Priority & Finding & Source lines \\
\midrule
\endfirsthead
\toprule ID & Priority & Finding & Source lines \\\midrule
\endhead
F01 & Critical & The candidate is checked against the principal root of
\code{problem^r}, not against \code{problem}. & 90--94, 121--134 \\
F02 & Critical & The custom factorizer merges nested exponents without a
branch condition. & 527--531 \\
F03 & Critical & Common-factor extraction under a noninteger outer power
assumes an invalid general product-power identity. & 538--543, 571--583 \\
F04 & High & Default \code{PossibleZeroQ} is used as an exact-equality
certificate. & 125--134 \\
F05 & High & Zero/invalid multipliers and an empty candidate set are not
safely rejected. & 107--134, 213--215 \\
F06 & High & The prime comparator compares the prime and exponent of its
\emph{first} argument. & 177--194 \\
F07 & High & \code{multipliercap = 1000} does not bound candidate count or
materialization cost. & 162--199 \\
F08 & Medium & Loading the file leaves \code{Print} unprotected; helper
symbols occupy the caller's context. & 1--4 and global definitions \\
F09 & High & A reported success does not establish reduced radical depth,
radical-only output, or final equality with the original input. & 120--140,
309--323, 434--440 \\
F10 & Medium & Input-domain checks are incomplete and occur after global
preprocessing in the wrapper. & 5--12, 79--108, 443--454 \\
F11 & Medium & Search heuristics have no completeness certificate; partial
factorization and seed normalization discard information. & 159--247,
396--420 \\
F12 & High & There is no global visited-multiplier set or overall work budget;
stack growth and expensive trials can dominate. & 251--430 \\
F13 & Medium & Positional, mutable state and implicit outcome tags make
invariants and error handling fragile. & 105--106, 253--266, 292--309 \\
F14 & Low/medium & Dead settings, a five-pass factorization bound, and the
meaning of \code{alllevels} can mislead maintenance and users. & 95--102,
162--165, 398--420, 520--549 \\
\bottomrule
\end{longtable}
\normalsize

\section{Architecture and public behavior}
\label{sec:architecture}
\subsection{Function map}
The twelve top-level definitions divide naturally into an orchestration layer,
a search layer, and structural/algebraic helpers.
\begin{center}\small
\begin{tabularx}{\textwidth}{@{}l r X@{}}
\toprule
Function & Lines & Role \\
\midrule
\code{Strad} & 1--3 & Calls the wrapper while dynamically suppressing
\code{Print}. \\
\code{DenestRadicals3} & 4--76 & Finds and groups radical tasks, then invokes
a chosen denester. \\
\code{denest11} & 78--441 & Builds multipliers, tests polynomial degrees and
GCDs, and manages a search stack. \\
\code{RationalizeDenominator} & 443--451 & Computes an algebraic denominator's
inverse from its minimal polynomial. \\
\code{AlgebraicQ} & 453--454 & Tests membership in \code{Algebraics}. \\
\code{nestdepths} & 456--465 & Annotates an expression tree with a selected
nesting count. \\
\code{nestdepthsort} & 468--479 & Groups those annotations by their nesting
count. \\
\code{maxnestdepth} & 482--497 & Extracts nodes with maximal selected nesting. \\
\code{ComplexitySort} & 500--504 & Orders candidates by leaf count, then
absolute value. \\
\code{Factorc} & 506--516 & Decodes the custom factor representation. \\
\code{fullfactor} & 518--550 & Produces that representation, recursively
processing sums and powers. \\
\code{combine} & 553--583 & Extracts bases common to the terms of a sum. \\
\bottomrule
\end{tabularx}
\end{center}
Two important local functions, \code{trymultiplier} and
\code{newmultipliers}, are defined inside \code{denest11}. The local
\code{SumQ} belongs to the wrapper's preprocessing rule.

\subsection{Calling conventions and their limitations}
The apparent interfaces are
\begin{lstlisting}
Strad[arguments...]
DenestRadicals3[expr, alllevels : False, func : denest11]
denest11[problem, forcedmultipliers : 0]
\end{lstlisting}
The displayed colons summarize defaults; they are not replacement definitions.
Only a literal \code{True} enables the \code{alllevels} branch. A value such as
\code{Automatic} or a misspelling is treated like \code{False}, not rejected.
The third argument is a callable denesting engine. A list supplied as the
second argument to \code{denest11} supplies an initial multiplier batch; a
non-list, including its default \code{0}, selects heuristic initialization.
Thus \code{0} and \code{{0}} have very different meanings.

The name \code{forcedmultipliers} should not be understood as a documented
hard whitelist. Later improvements can enqueue generated batches even when
the initial list was supplied by the caller. An empty supplied list yields no
trials and leaves the problem unchanged. These are source-derived interface
observations, not evidence that such behavior was intentional.

\subsection{Preprocessing and task marking}
The wrapper first evaluates \code{Factorc[expr]} \emph{before} applying the
algebraicity predicate. It then recognizes expressions of the form
\begin{lstlisting}
radicand_?AlgebraicQ^pow_Rational
\end{lstlisting}
and temporarily replaces them by a local marker \code{f[radicand,pow]}
(\src{11--40}). For each matched radicand it:
\begin{enumerate}
\item obtains a common coefficient using \code{FactorTermsList};
\item factors the remaining part using the custom factorizer, and decomposes
rational or Gaussian-rational coefficient factors;
\item combines factors that contain a sum or an explicit complex number with
nonzero real part;
\item moves signs into a single sign accumulator, rationalizes denominators,
and simplifies the reconstructed product.
\end{enumerate}
The \code{SumQ} name is broader than ``is a sum'': its test searches all levels
for \code{Plus} or \code{Complex[Except[0], _]}. It is a heuristic classifier,
not a mathematical characterization of additive structure.

The sign-flip loop deserves careful interpretation. Replacing a factor $a_i$
by $-a_i$ and also multiplying an accumulated sign by $-1$ preserves the
\emph{product}. Consequently a numerical misclassification by itself does not
prove that this loop changes the input value. The serious branch defects
occur elsewhere, where fractional powers are reinterpreted or factored.

\subsection{Why tasks are combined}
Before calling the denester, the wrapper combines two marked factors when both
radicands contain a \code{Power} node (\src{43--47}). It then changes
\code{f[rad,pow]} to \code{f[rad^pow]}. A second rule combines one-argument
markers when both enclosed numerical values have nonzero real and imaginary
parts at a requested precision of ten digits (\src{49--54}).

These rules do not themselves assert $\sqrt{ab}=\sqrt a\sqrt b$. They form a
marker around an already evaluated product. That grouping can be useful
because a product may admit a joint simplification. It is nevertheless
important: a product of principal roots need not lie in the sector of a
principal root of the inferred degree. The grouped product can therefore
violate an assumption that \code{denest11} silently makes; see
Section~\ref{sec:branch}.

\subsection{The two meanings of traversal}
With \code{alllevels =!= True}, the wrapper uses \code{ReplaceAll}, gathers
distinct marked tasks, applies the chosen function once to each task, and
simplifies the result. With \code{alllevels === True}, repeated marking exposes
nested tasks, and a small dependency-resolution loop processes markers whose
arguments contain no unresolved markers. Resolved tasks become replacement
rules and are substituted into the remaining tasks
(\src{55--73}).

This is approximately an innermost-first traversal of a dependency graph. It
is \emph{not} a proof that every radical newly introduced by the denester is
reprocessed until no further improvement is possible. In particular, a fixed
point of the marker-replacement rules is not the same thing as a fixed point
of the complete denesting procedure. Wolfram's \code{ReplaceRepeated} itself
has a documented iteration limit by default; that built-in limit does not
provide a useful whole-program search budget \cite{replacerepeated}.

\section{The mathematical engine}
\label{sec:engine}
\subsection{From an expression to a binomial equation}
Let $\beta$ denote the original problem presented to the core. The program
chooses a positive integer $r$, called \code{reductionpower}, as the least
common multiple of denominators attached to selected maximal-depth rational
powers. It then sets
\[
 \code{outerpower}=\frac1r,\qquad R=\code{radicand}=\beta^r.
\]
A further product of denominator LCMs, called \code{extrapower}, is computed
and printed, but is never used to change the actual search
(\src{90--102}). Thus the ``extra root'' diagnostic does not mean that an
extra-root algorithm is implemented.

For a proposed nonzero algebraic multiplier $m$, the trial routine forms
\[
 \alpha=(mR)^{1/r},\qquad p_m(X)=\operatorname{MinPoly}_{\Q}(\alpha),
 \qquad d_m=\deg p_m.
\]
Here and throughout, a fractional power denotes the principal complex power.
The one-argument form \code{MinimalPolynomial[alpha]} in the attachment is
valid: it returns a pure function, so applying the returned value to
\code{arg1} gives a polynomial \cite{minimalpolynomial}.

When its degree-based screening condition is satisfied, the program computes
\begin{equation}
 g_m(X)=\gcd\bigl(p_m(X),\,X^r-mR\bigr)
 \quad\text{over an algebraic coefficient field}.\label{eq:gcd}
\end{equation}
The option \code{Extension -> Automatic} matters: it tells
\code{PolynomialGCD} to include algebraic numbers appearing in the
coefficients, rather than treating all displayed radicals as independent
indeterminates \cite{polynomialgcd}. The effective field can depend on the
representations of those coefficients. It need not be described merely as
$\Q(\sqrt{c})$ in more complicated examples.

\subsection{What a proper GCD proves}
\begin{proposition}\label{prop:gcd}
Assume $r\ge2$, $m\ne0$, and $R$ are exact algebraic numbers. Let $K$ be a
characteristic-zero field containing the coefficients used in
\eqref{eq:gcd}. Then $g_m$ has positive degree. Every root of $g_m$ is a root
of $X^r-mR$. If $\deg g_m<r$, the GCD gives an equation of lower degree for
at least one admissible scaled root.
\end{proposition}
\begin{proof}
The number $\alpha=(mR)^{1/r}$ is a common root of $p_m$ and $X^r-mR$ in an
algebraic closure of $K$. If the polynomials were relatively prime over $K$,
Bezout's identity would give $A p_m+B(X^r-mR)=1$ for suitable polynomials
$A,B\in K[X]$. Evaluation at $\alpha$ would give $0=1$, a contradiction.
Thus the GCD is nonconstant. Because it divides the binomial, each of its
roots satisfies the binomial. The final assertion is precisely the
inequality $\deg g_m<r$.
\end{proof}
The proposition justifies the central computation. It does not say that
solving the smaller equation produces a simpler \emph{radical expression}.
Its coefficients may already be complicated, and the solver may return
\code{Root} objects or introduce additional radicals. The program conflates
these different notions of progress.

\subsection{Why the roots of unity are appropriate}
For each root $a$ returned by \code{Solve[g_m == 0, X]}, the program generates
\begin{equation}
 \frac{a}{m^{1/r}}\zeta_r^k,
 \qquad \zeta_r=e^{2\pi i/r},\quad 0\le k<r.\label{eq:candidates}
\end{equation}
This is a sound way to handle the phase introduced by scaling, provided the
right original target is retained.
\begin{theorem}\label{thm:phase}
Let $m\ne0$ and $\beta\ne0$ be complex numbers, let $r\ge2$ be an integer,
and put $R=\beta^r$. For any $a$ with $a^r=mR$, the set in
\eqref{eq:candidates} is exactly the set of all $r$ roots of $X^r=R$.
In particular, it contains $\beta$.
\end{theorem}
\begin{proof}
Set $\mu=m^{1/r}$. Whatever principal phase is selected, $\mu^r=m$ and
$\mu\ne0$. Hence $(a/\mu)^r=R$. Multiplication by $\zeta_r^k$ preserves this
equation. These $r$ numbers are distinct because $a/\mu\ne0$ and the roots of
unity are distinct. A degree-$r$ polynomial has at most $r$ roots, so they are
all the roots. Since $\beta^r=R$, $\beta$ occurs among them.
\end{proof}
No identity $(mR)^{1/r}=m^{1/r}R^{1/r}$ is needed in this proof. That identity
is false in general. The phase enumeration is therefore a strength of the
program, not a defect. Its defect is that it selects the branch using
$R^{1/r}$ instead of $\beta$.

If the GCD has $d$ roots, using every root and every phase produces $dr$
formal candidates. For branch coverage, one GCD root already suffices by the
theorem. Additional roots can still be useful as alternative symbolic
representations, but their role should be distinguished from obtaining new
algebraic values.

\subsection{A direct successful case}
Take $R=3+2\sqrt2$ and $r=2$. Since $1+\sqrt2>0$ and
$(1+\sqrt2)^2=R$, the desired principal root is $1+\sqrt2$. With $m=1$,
\[
 p_1(X)=X^2-2X-1,
 \qquad
 \gcd\bigl(p_1(X),X^2-3-2\sqrt2\bigr)=X-1-\sqrt2
\]
over $\Q(\sqrt2)$. The GCD has degree one, smaller than two, and the candidate
selection can recover the answer immediately. These polynomial identities
were independently checked exactly.

\subsection{A case where a multiplier helps}
Let
\[
 R=5+2\sqrt6=(\sqrt2+\sqrt3)^2.
\]
The unscaled principal root has minimal polynomial
\[
 p_1(X)=X^4-10X^2+1,
\]
and over $\Q(\sqrt6)$ the GCD with $X^2-R$ is the entire binomial. Thus this
trial does not give a lower-degree equation. With $m=2$, however,
\[
 (2+\sqrt6)^2=2R,
 \qquad p_2(X)=X^2-4X-2,
\]
and the GCD becomes $X-2-\sqrt6$. Dividing the resulting root by $\sqrt2$
gives $\sqrt2+\sqrt3$.

The discriminant of $p_1$ is
\[
 \Disc(p_1)=147456=2^{14}3^2.
\]
Consequently the discriminant-based generator has access to the primes 2 and
3 and can propose $m=2$. This explains the heuristic's motivation without
claiming a particular actual trial order: the comparator defect can affect
that order.

\section{Multiplier search and its state machine}
\label{sec:search}
\subsection{Initialization from visible terms}
Unless a list of multipliers is supplied, the core expands the rationalized
radicand, strips visible rational coefficients from its terms, discards pure
rational terms, and replaces certain complex terms by a signed imaginary
unit (\src{224--247}). It then constructs complementary powers of visible
radicals. For a factor displayed as $b^{a/q}$, a typical calculation is
\[
 \left(b^{a/q}\right)^{(q-a)/a}.
\]
In a positive-real setting this is a way of suggesting a complementary factor
$b^{(q-a)/q}$. In the complex setting it is only a candidate-generation
heuristic, not a universally valid exponent simplification. Because these
expressions are \emph{multipliers to be tried}, changing their value is not by
itself an incorrect rewrite of the original problem. It may instead make the
search less effective.

The initialization also assumes that \code{List @@ Expand[...]} is a list of
summands. That is valid when the head is \code{Plus}. If the expanded expression
has head \code{Times} or \code{Power}, the same operation exposes factors or
arguments instead. For an atom it does not magically produce a singleton
list. A reliable implementation should explicitly use
\begin{lstlisting}
terms = If[Head[expanded] === Plus, List @@ expanded, {expanded}];
\end{lstlisting}
The original nested-power precheck does not prove this missing summand
invariant.

\subsection{Degree history and outcome formats}
A trial starts with a record of the form
\[
 \{\texttt{False},\{m,p_m,d_m\}\}.
\]
It may instead return
\[
 \{\texttt{True},\{m,p_m,d_m\}\}
 \quad\text{or}\quad
 \{\texttt{True},\{m,p_m,d_m\},\text{answer}\}.
\]
The first Boolean means ``search progress,'' not ``a correct denesting was
found.'' In particular, a larger degree may be recorded as progress when it
reveals that a previously observed smaller degree was significant
(\src{147--151}).

The history starts with the sentinel $\{0,\infty,\infty\}$ and is maintained
in a degree-based order. References such as \code{history[[-1]]} therefore mean
a currently best degree record, not necessarily the most recent chronological
attempt. The code's use of GCDs of degrees can guide exploration, but degree
arithmetic alone is not a certificate of field membership or denestability.

\subsection{Generating batches from discriminants}
For comparable observed degrees, \code{newmultipliers} partially factors the
discriminant of a minimal polynomial, removes the sign factor, discards factors
that fail \code{PrimeQ}, and tries products of selected primes with exponents
between zero and $r-1$. The command
\code{FactorInteger[n, Automatic]} deliberately extracts only factors that
are easy to find; an unresolved composite cofactor can remain
\cite{factorinteger}. Discarding that cofactor loses access to its unknown
prime divisors.

There is an important qualification. Primes removed from the
\emph{combination prefix} by the ratio or length heuristic are added back as
\emph{individual} candidates through
\code{Union[Divisors[...], primes]}. They are not all discarded outright.
What is lost is their participation in most products, together with every
prime hidden inside a discarded unresolved composite.

When the GCD of two observed degrees is smaller than either degree, the other
branch simply proposes the product of the two recorded multipliers
(\src{201--207}). The attachment supplies no theorem making this product
sufficient. As a general warning against reasoning from degree GCDs alone,
$\sqrt2$ has degree 2, $\sqrt[3]3$ has degree 3, and their product has degree 6,
not degree $\gcd(2,3)=1$. This example is not asserted to be an execution trace
of that particular branch; it illustrates why the heuristic needs a separate
justification.

\subsection{Why the discriminant is suggestive, not conclusive}
For an algebraic integer $\alpha$ generating a number field $K$, the familiar
relation
\begin{equation}
 \Disc(p_\alpha)=\Disc(K)
 [\mathcal O_K:\Z[\alpha]]^2\label{eq:disc}
\end{equation}
shows that a polynomial discriminant includes both field information and
information about the chosen generator. To see the square, compare a power
basis with an integral basis in the trace-pairing matrix: a basis change by a
matrix $A$ transforms the matrix to $A^{\mathsf T}TA$, whose determinant is
$(\det A)^2\det T$. The absolute determinant of this lattice change is the
index in \eqref{eq:disc}.

For example, $\alpha=3\sqrt2$ has minimal polynomial $X^2-18$ and discriminant
$72=8\cdot3^2$, whereas the usual integral basis of $\Q(\sqrt2)$ has
discriminant 8. Thus the generator introduces the additional prime 3. For
nonintegral generators one cannot apply \eqref{eq:disc} to a primitive,
nonmonic polynomial without further qualification. In either setting, a list
of discriminant primes is not automatically a complete or minimal list of
useful denesting multipliers.

\subsection{Batch scheduling}
The stack stores either a concrete candidate list or a delayed call to
\code{newmultipliers}. Some batches are marked ``do full set.'' Other batches
may be interrupted when progress creates a more promising batch, with the
untried remainder inserted separately (\src{251--430}). At the end of the
initial heuristic pass, additional products of promising seeds are formed by
\code{Distribute} and further exponent normalization.

This is a reasonable exploratory scheduling idea, but the representation is
fragile. A single position may hold an unevaluated multiplier, an evaluated
pair $\{m,d\}$, or a delayed generator specification, depending on the phase of
the search. An answer is recognized using
\code{ListQ[attempt[[-1]]] =!= True}, rather than an explicit outcome tag.
That works only while the undocumented invariant ``an answer is never a
list'' and the assumed attempt shape both hold.

\section{The helpers: correct mathematics and unsafe reconstruction}
\label{sec:helpers}
\subsection{Denominator rationalization is algebraically correct}
Let $d\ne0$ be algebraic and let
\[
 p(X)=a_0+a_1X+\cdots+a_sX^s
\]
be its minimal polynomial over $\Q$, normalized in any nonzero rational
multiple. Since $d\ne0$, $a_0\ne0$. Dividing the identity $p(d)=0$ by $d$
gives
\begin{equation}
 \frac1d=-\frac{a_1+a_2d+\cdots+a_sd^{s-1}}{a_0}.
 \label{eq:inverse}
\end{equation}
This is exactly the identity implemented in \src{443--450}.

More explicitly, divide $p(X)$ by $X-d$ to obtain $p(X)=(X-d)q(X)$, because
the remainder is $p(d)=0$. Evaluating at zero gives
\[
 p(0)=-d\,q(0),\qquad \frac{q(0)}{-p(0)}=\frac1d.
\]
The source multiplies this last quotient by the original numerator. Its minus
sign is correct, and a monic minimal polynomial is not required. The name
\code{conjugate} is somewhat misleading: it stores a minimal polynomial, not
one selected complex conjugate.

For $d=1+\sqrt2$, the polynomial is $X^2-2X-1$, giving
$d^{-1}=\sqrt2-1$. For $d=1+i$, it is $X^2-2X+2$, giving
$d^{-1}=(1-i)/2$. Both, along with three further examples, were checked by
exact polynomial division in the accompanying Python script.

The limitations are contractual and computational. The helper does not guard
against a zero or nonalgebraic denominator, does not constrain the cost of
\code{MinimalPolynomial}, and does not first apply \code{Together}. A sum of
several fractions can have top-level denominator 1 even though denominators
remain in its terms. Some callers explicitly use \code{Together}, but the
initialization at \src{224--229} does not. Expanding the inverse in
\eqref{eq:inverse} can also make an expression substantially larger.

\subsection{What the nesting helpers measure}
For a predicate $P$ on expression nodes, the intended annotation follows the
recurrence
\[
 h_P(E)=\mathbf1_{P(E)}+
 \max\bigl(\{h_P(c):c\text{ is an argument of }E\}\cup\{0\}\bigr).
\]
The source stores the original expression, its child annotations, and its
height in nested lists. With $P(E)$ equal to ``$E$ is a rational power,'' this
is a \emph{syntactic radical height}. It is not an invariant of the algebraic
number. The same number can have several heights, and a \code{Root} object
may conceal substantial algebraic complexity while containing no displayed
rational power at all.

\code{nestdepthsort} flattens this ad hoc annotation into height groups;
\code{maxnestdepth} repeatedly prunes entries below the maximum height. The
code is compact but depends on list shapes, \code{Sequence} splicing, and the
distinction between expression data and structural lists. General list-valued
input is not given a reliable independent contract. A cleaner implementation
would return explicit records such as $\{\text{position},\text{height}\}$ and
would handle the empty set of matching powers deliberately.

The early test in \code{denest11} is also broader than the later height test:
it recognizes an outer \code{Power} whose base contains any \code{Power},
whereas the denominator computation selects only \code{_Rational} exponents.
For example, a symbolic expression $(1+x^2)^3$ can satisfy the former test
without containing a rational-exponent radical. The code does not establish
that the resulting \code{reductionpower} is an integer at least two before
using it. The trial helper is defined only for an outer exponent matching
\code{Rational[1, _Integer]}, so an integer outer exponent 1 is not an
interchangeable case.

\subsection{The custom factor representation}
\code{fullfactor} starts from \code{FactorList}, whose entries are
factor--exponent pairs \cite{factorlist}. It recursively factors nonprime
integer and rational bases, pulls exponents through displayed powers, and
processes sums term by term using \code{combine}. The latter merges identical
bases within each term, finds bases common to all terms, takes the minimum
exponent for each common base, and subtracts it from every term.

\code{Factorc} then reconstructs ordinary expressions through several
shape-based rules, using the listability of \code{Times} and \code{Plus} and
finally extracting the first item of a flattened result. This is not merely a
cosmetic replacement for the built-in factorizer: it introduces new algebraic
rewrite rules. Those rules require branch conditions that the implementation
does not check.

\section{Correctness defects and counterexamples}
\label{sec:correctness}
\subsection{F01: the original branch is discarded}
\label{sec:branch}
The central incorrect inference is
\[
 \beta\longmapsto R=\beta^r
 \quad\text{followed by the assumption}\quad R^{1/r}=\beta.
\]
For a principal root, this is false. Write $\beta=\rho e^{i\theta}$ with
$\rho>0$ and $-\pi<\theta\le\pi$. There is an integer $k$ such that
\[
 \Arg(\beta^r)=r\theta-2\pi k\in(-\pi,\pi].
\]
Therefore
\begin{equation}
 (\beta^r)^{1/r}=\beta e^{-2\pi i k/r}.
 \label{eq:branchloss}
\end{equation}
The reconstructed root equals $\beta$ precisely when $\beta$ lies in the
principal root sector $-\pi/r<\Arg\beta\le\pi/r$. Neither the general core
interface nor the wrapper's product-grouping rules establish this condition.
Principal-value semantics are the documented behavior of \code{Power}
\cite{power}.

The source's final filter is
\begin{lstlisting}
Select[possibilities, PossibleZeroQ[# - radicand^outerpower] &]
\end{lstlisting}
The expression after the minus sign should be the \emph{original problem},
not its reconstructed principal root. Merely replacing \code{PossibleZeroQ}
with an exact method would not fix this wrong target.

\subsubsection{An exact negative-real counterexample}
Take the core problem
\[
 \beta=-\sqrt{3+2\sqrt2}=-(1+\sqrt2),\qquad r=2.
\]
The trial path, with multiplier $m=1$, has
\[
 R=3+2\sqrt2,\quad p(X)=X^2-2X-1,\quad
 g(X)=X-1-\sqrt2.
\]
Its phase candidates are $1+\sqrt2$ and $-1-\sqrt2$. The code compares them to
$R^{1/2}=1+\sqrt2$ and selects the \emph{positive} candidate. The original
negative candidate is the one it should preserve. This is an exact sign
error, not a numerical-precision issue.

The corresponding regression call is
\begin{lstlisting}
denest11[-Sqrt[3 + 2 Sqrt[2]], {1}]
\end{lstlisting}
The algebraic trial path above is proved and independently checked. The call
itself remains an unexecuted Wolfram regression probe; automatic evaluation
can affect whether a written nested expression survives to enter that path.
That qualification does not make the normalization identity valid.

\subsubsection{A complex product relevant to the wrapper}
The defect is not restricted to a caller deliberately putting a minus sign
in front of a core problem. Define
\[
 u=1+i\sqrt2,\qquad v=1+i\sqrt3.
\]
Since both have positive real part,
\[
 \sqrt{-1+2i\sqrt2}=u,\qquad
 \sqrt{-2+2i\sqrt3}=v.
\]
Their product is
\begin{equation}
 \beta=uv=1-\sqrt6+i(\sqrt2+\sqrt3),\label{eq:complexbeta}
\end{equation}
which has negative real part and positive imaginary part. Hence
$\sqrt{\beta^2}=-\beta$. Explicitly,
\[
 R=\beta^2=2-4\sqrt6-i(4\sqrt2+2\sqrt3).
\]
The reconstructed principal root $\gamma=-\beta$ has minimal polynomial
\[
 p_\gamma(X)=X^4+4X^3+4X^2+48X+144,
\]
and over a coefficient field containing the displayed radicals,
\[
 \gcd(p_\gamma(X),X^2-R)
   =X-\sqrt6+1+i\sqrt2+i\sqrt3=X-\gamma.
\]
The source again filters for $\gamma$, not for $\beta$.

This example lies in exactly the kind of product class that the wrapper's
marker-combination rules create. The preceding calculation establishes the
core branch defect for the grouped algebraic value. The public-wrapper probe
in the archive tests the additional interaction with \code{Factorc} and
automatic evaluation; no unobserved wrapper transcript is claimed.

\subsubsection{Required repair}
Pass the original target explicitly into every trial, and select by a
certified equality to that target. The roots-of-unity enumeration already
contains the correct phase under the hypotheses of Theorem~\ref{thm:phase}.
The repair does not require guessing a sign from approximate real parts.
For an exact algebraic target, a minimal acceptance pattern is
\begin{lstlisting}
Select[candidates, TrueQ[RootReduce[# - original] === 0] &]
\end{lstlisting}
This must be accompanied by input, failure, and empty-list handling.

\subsection{F02: unguarded multiplication of nested exponents}
At \src{527--531}, a factor record whose base is itself a power is changed
from the mathematical shape
\[
 \{a^b,c\}\quad\text{to}\quad\{a,bc\}.
\]
That represents the rewrite $(a^b)^c\mapsto a^{bc}$ without a branch test.
The rewrite is not valid for general principal powers. For example,
\[
 \sqrt{(-2)^2}=2\ne-2.
\]
For a symbolic input, the normalization can move through the conceptual
records
\[
 \{\sqrt{x^2},1\}\longmapsto\{x^2,1/2\}
 \longmapsto\{x,1\}.
\]
The result $x$ does not equal $\sqrt{x^2}$ on the real line: the correct
real-valued expression is $|x|$. The regression suite therefore checks
\code{Factorc[Sqrt[t^2]]} and only \emph{afterward} substitutes $t=-2$.
Substituting before the call could allow ordinary evaluation to hide the
faulty symbolic rule.

The outer wrapper makes this especially important. Its first operation is
\code{Factorc[expr]}, before the \code{AlgebraicQ} pattern restriction. Thus a
user can pass symbolic input that is altered unsafely rather than rejected or
left unchanged. Even if the intended public domain is only explicit algebraic
constants, the implementation does not enforce that domain before applying
the problematic rewrite.

A safe repair needs a theorem for the permitted case: for instance, a
positive real base with suitable real exponents, or an integer outer
exponent. Otherwise the nesting must be retained. A blanket call to
\code{PowerExpand} would repeat the same fundamental mistake rather than
repair it.

\subsection{F03: factoring a common base under a root}
The common-factor logic of \code{combine} is valid for a sum before an outer
power is taken. It becomes unsafe when a common factor is assigned exponent
\code{factor*pow} outside the residual sum. The resulting identity is
\[
 (cS)^q\stackrel{?}{=}c^qS^q,
\]
which is not generally true for noninteger $q$.

For a concrete symbolic example, factoring the common $x$ in $x+x^2$ suggests
\[
 \sqrt{x+x^2}\stackrel{?}{=}\sqrt{x}\sqrt{1+x}.
\]
At $x=-2$, the left side is $\sqrt2$, whereas the right side is
$(i\sqrt2)i=-\sqrt2$. This is the behavior targeted by the
\code{Factorc[Sqrt[t + t^2]]} substitution probe. The formal common-base
records for the two terms are $\{\{x,1\}\}$ and $\{\{x,2\}\}$, so the source
extracts exponent 1 before multiplying it by the outer exponent $1/2$.

The phase issue also occurs for explicit algebraic complex numbers. Put
\[
 B=\sqrt2+(-1)^{1/3}
   =\sqrt2+\tfrac12+\tfrac{\sqrt3}{2}i.
\]
Then $0<\Arg B<\pi/2$, so
\[
 \sqrt{-B}=-i\sqrt B,
 \qquad \sqrt{-1}\sqrt B=i\sqrt B.
\]
Extracting a common $-1$ from the terms under this root reverses the sign.
This proves that restricting coefficients to algebraic numbers, without a
positivity or branch condition, does not justify the rule.

The repair is to keep factored products \emph{inside} the noninteger power
unless a valid branch condition is established. Positive real factors can be
extracted safely for real rational exponents, subject to ordinary nonzero
domain conditions for negative exponents. General complex factors require
explicit phase accounting or exact comparison against the unsplit original.

\subsection{F04: a heuristic zero test is not a certificate}
Default \code{PossibleZeroQ} uses symbolic and numerical heuristics. The
official documentation explicitly describes possible incorrect positive
answers, including examples involving explicit algebraic quantities perturbed
by a very small nonzero rational number. It documents
\code{Method -> "ExactAlgebraics"} for guaranteed decisions on explicit
algebraic numbers \cite{possiblezero}.

Therefore the current filter cannot be advertised as an exact correctness
certificate. A false positive may admit a wrong candidate; failure to recognize
a true zero can also leave no acceptable candidate. Increasing a numerical
precision elsewhere in the code does not turn this test into a proof.

For the intended exact algebraic domain, a conservative certificate is
\code{RootReduce[candidate - original] === 0}. The alternative exact mode of
\code{PossibleZeroQ} is also relevant, but it must be applied to the
\emph{correct original target}. A timeout, unevaluated result, or unsupported
input must mean ``not certified,'' not ``approximately equal, therefore
accepted.'' This is a local algebraic equality decision, not a claim that
arbitrary symbolic equality is decidable.

\subsection{F05: zero multipliers and unchecked indexing}
The multiplier interface accepts an arbitrary list without checking that its
elements are nonzero exact algebraic numbers. Consider a quadratic trial with
$m=0$. Then
\[
 \alpha=0,\qquad p(X)=X,\qquad
 \gcd(X,X^2)=X.
\]
The GCD passes the source's lower-degree test. Solving gives $a=0$, after which
\eqref{eq:candidates} attempts to divide by $0^{1/2}$, producing undefined
arithmetic. The following supplied-input probe therefore has a definite
invalid path:
\begin{lstlisting}
denest11[Sqrt[3 + 2 Sqrt[2]], {0}]
\end{lstlisting}
The precise subsequent messages and returned symbolic debris depend on kernel
error propagation, and are not claimed here as observed output.

Independently of this trigger, the accepted candidates are stored in
\code{result}, after which the source accesses \code{result[[1]]} without first
checking whether the list is empty (\src{125--134}). An empty list may follow
invalid arithmetic, an unexpected solver result, a failed equality test, or
an interrupted algebraic computation. A proper guard is mandatory:
\begin{lstlisting}
If[valid === {},
  Return[Failure["NoCertifiedCandidate", <||>]]
];
\end{lstlisting}
Every proposed multiplier must be rejected before division unless it is a
supported, certified nonzero value. The problem $\beta=0$ should be handled
separately as a trivial case.

\subsection{F09: the word ``success'' overstates the guarantee}
The source declares a denesting when the GCD degree drops and an accepted
candidate is returned. It does not enforce any of the following independent
properties:
\begin{enumerate}
\item equality to the original expression before preprocessing;
\item a decrease in explicit radical nesting or another chosen complexity
measure;
\item output in a specified radical-only expression language.
\end{enumerate}
A lower-degree polynomial equation need not yield fewer nested radicals.
Furthermore, \code{Solve} can return \code{Root} objects; for example, its
documented default treatment of numerical cubic equations can use such
objects \cite{solve}. Hiding an algebraic number inside \code{Root} is not the
same task as denesting its radical representation.

The final comparison \code{solution =!= problem} is a structural comparison,
not a semantic equality test or a simplicity ordering. Multiple retained
candidates are likewise not necessarily different solutions: ordinary
\code{DeleteDuplicates} can leave distinct expressions for the same algebraic
number. If every candidate is exactly equal to the target, they are all the
same \emph{value}. Select the best representation by an explicit metric,
then deduplicate with exact algebraic equality if needed.

\section{Defects in search control, interfaces, and software engineering}
\label{sec:software}
\subsection{F06: the prime comparator uses the wrong slot}
The comparator in \src{181} is
\begin{lstlisting}
#1[[1]] <= #[[2]] &
\end{lstlisting}
In a pure function, \code{#} means \code{#1}, not \code{#2}
\cite{function}. If the first factorization pair is $\{p,e\}$ and the second
is $\{q,f\}$, the predicate evaluates $p\le e$. It never examines $q$ or $f$.
For example,
\[
 \operatorname{bad}(\{2,1\},\{3,1\})=\texttt{False},\qquad
 \operatorname{bad}(\{3,1\},\{2,1\})=\texttt{False}.
\]
It therefore does not order the primes. With a first pair $\{2,18\}$, the
same predicate is true no matter what the second pair is.

This affects more than presentation: subsequent prime-ratio checks and prefix
selection rely on ascending prime order. The most direct repair is
\begin{lstlisting}
primes = SortBy[Select[primes, PrimeQ[First[#]] &], First];
\end{lstlisting}
Alternatively, the intended binary comparator is
\code{#1[[1]] <= #2[[1]] &}. The other uses of \code{<=} in sorting should not
be condemned merely because another programming language requires a strict
weak comparator; Wolfram's sorting interface explicitly supports Boolean
ordering predicates \cite{sort}. The demonstrated defect is the wrong slot.

\subsection{F08: global protection and context pollution}
The second line of the attachment executes \code{Unprotect[Print]}. No later
line restores that attribute. The dynamic binding in
\begin{lstlisting}
Block[{Print = (Null &)}, DenestRadicals3[eaf]]
\end{lstlisting}
restores the temporarily localized \emph{value} on leaving the block, but it
does not undo the separate earlier removal of \code{Protected}. The
\code{Unprotect} operation changes attributes; \code{Block} localizes values
\cite{unprotect,block}. Thus merely loading this file changes the global
session's protection state.

A preferable design routes diagnostics through a private logger or a
\code{Verbose} option and never alters \code{System`Print}. Suppressing
\code{Print} is also not the same as suppressing kernel messages, and
arguments to \code{Strad} ordinarily evaluate before its body is entered.

The absence of \code{BeginPackage} and a private context also matters. Loading
the file clears and redefines generic names such as \code{combine},
\code{Factorc}, and \code{AlgebraicQ} in the active context. This is normal for
an isolated notebook experiment but risky for a reusable package. In
addition, dynamic \code{Block} localization does not create fresh marker
symbols. A caller's attributes on a symbol such as \code{f} can matter.
\code{Module} creates fresh symbols and is a better default for private
implementation state \cite{module}. This is not a claim that every use of
\code{Block} in the source is wrong.

\subsection{F10: the input domain is not enforced}
The mathematical core needs at least explicit exact algebraic inputs,
nonzero multipliers, and an integer $r\ge2$. The public definitions accept
much more. They do not provide consistent behavior for approximate numbers,
symbolic parameters, transcendental constants, list-valued inputs, invalid
options, or malformed trial outcomes.

\code{AlgebraicQ[num_] := Element[num, Algebraics]} is not automatically a
bug just because \code{Element} may remain symbolic. A \code{PatternTest}
accepts only literal \code{True}; an undecided predicate fails the match
\cite{element,patterntest}. The genuine problems are that the predicate is
not a full API contract and is applied after the global custom factorizer.
The direct core has still fewer guards.

For a deliberately numerical algebraic interface, an initial predicate can
combine \code{NumericQ}, infinite \code{Precision}, and a literal algebraic
membership result. Using precision also catches inexact complex atoms; merely
searching for visible \code{_Real} subexpressions is not a complete precision
policy. A separate symbolic interface would need explicit assumptions and
assumption-preserving transformations. \code{ComplexExpand} assumes unlisted
variables are real \cite{complexexpand}; that is another reason not to pass
unrestricted symbolic expressions through this numeric preprocessing pipeline.

\subsection{F11--F13: incompleteness, repeated work, and fragile state}
The search is not proved exhaustive. It screens trials by minimal-polynomial
degree, uses visible-term guesses, partially factors discriminants, limits
prime combinations, and normalizes away some seed factors. A failure to find
an improvement proves only that the selected search did not find one. It is
not a certificate that no denesting exists.

The stack also has no global set of previously tried exact multipliers.
Different generated batches can contain the same multiplier, and structurally
different expressions can denote the same multiplier. A local
\code{DeleteDuplicates} or \code{Union} does not provide global semantic
memoization. Repeating \code{MinimalPolynomial} and \code{PolynomialGCD} for
these values can be very expensive.

The state representation compounds the problem. Positional indexing such as
\code{multipliers[[1, 2, {1, -1}]]} assumes several nested layouts at once.
The same stack field changes interpretation as candidates are evaluated.
The Boolean named \code{stackchange} is set when progress is noticed, not
solely at a physical stack insertion. The logic is difficult to audit and
requires tests of interruption and resumption invariants. No particular
infinite loop or universal stack-corruption trace is asserted here; the
source simply lacks a clear, enforced state contract.

Use named records with distinct variants: an unevaluated candidate, an
attempt result, a pending generator, and a queue batch. Suggested trial
outcomes are \code{"NoProgress"}, \code{"VerifiedCandidate"},
\code{"InvalidInput"}, \code{"TimedOut"}, and \code{"InternalFailure"}.
``Found an expression'' and ``made a search-scheduling improvement'' should
not share a Boolean flag.

\subsection{F14: dead settings and ambiguous maintenance cues}
Several details are demonstrable data-flow shortcomings rather than
mathematical counterexamples. The computed \code{extrapower} is only printed.
\code{maxprimes = 14} is never read. The local \code{best} assigned in the
post-initial-pass block is unused, whereas \code{worst} is used. Some local
names and comments describe abandoned variants, and substantial obsolete
code remains commented out.

The \code{j < 5} bound in \code{fullfactor} is a real implementation limit.
It prevents an unlimited local normalization loop, but there is no fixed-point
check or status reporting when the fifth pass still changes the representation.
It should not be described as a proof of complete factoring. Conversely, the
existence of this bound means it would be inaccurate to claim that every loop
in the program is unbounded.

\section{Cost model and the false multiplier cap}
\label{sec:cost}
\subsection{F07: exact candidate counts}
Let $t$ distinct primes be selected for combination and suppose $r\ge2$.
The integer fed to \code{Divisors} is
\[
 N=\left(\prod_{j=1}^{t}p_j\right)^{r-1}.
\]
Every divisor has the unique form
\[
 \prod_{j=1}^{t}p_j^{e_j},\qquad 0\le e_j\le r-1,
\]
so $N$ has exactly $r^t$ divisors. Dropping the unit leaves $r^t-1$ candidates.
If there are $s$ detected primes in all, with the other $s-t$ reintroduced only
as singletons, the actual union has $r^t-1+(s-t)$ candidates before scaling.
For a nonzero fixed multiplier, scaling does not identify distinct values.

When enough primes survive the prefix heuristics, the source uses
\[
 t=\max\left(5,\left\lceil\frac{\log1000}{\log r}\right\rceil\right).
\]
The ceiling, the minimum of five selected primes, and the later singleton
union are all incompatible with treating 1000 as a hard cap.
\begin{center}
\begin{tabular}{@{}rrr@{}}
\toprule
Root degree $r$ & Selected primes $t$ & $r^t-1$ \\
\midrule
2 & 10 & 1,023 \\
3 & 7 & 2,186 \\
4 & 5 & 1,023 \\
5 & 5 & 3,124 \\
6 & 5 & 7,775 \\
10 & 5 & 99,999 \\
100 & 5 & 9,999,999,999 \\
\bottomrule
\end{tabular}
\end{center}
These are conditional counts assuming the indicated number of primes is
available. They are not measured timings or assertions that every input
produces that many candidates. The final row nevertheless shows the scale of
the missing safeguard.

A cap must be enforced \emph{before} candidate materialization. Computing all
\code{Divisors} and then taking the first 1000 does not prevent the allocation.
A bounded iterator over exponent vectors, preferably ordered by a useful
cost, is a more appropriate generator.

\subsection{Other sources of growth}
The post-initial-pass \code{Distribute} operation forms products by taking or
not taking selected seed multipliers. With $k$ such choices, the unpruned
product family can have $2^k$ members. Subsequent normalization may collapse
many of them, but it does not remove the cost of forming the intermediate
family.

For an individual trial, the important parameters are not only the number of
candidates but also the root degree $r$, the degrees of the algebraic
coefficient fields, the bit sizes of coefficients, and the expression size of
multipliers. Minimal polynomials of combinations of independent algebraic
numbers can have degrees approaching the product of the constituent degrees.
Expanding $\beta^r$ and rationalized denominators can sharply increase both
coefficient height and expression size. Solving a proper GCD of degree $d$
and forming every phase yields up to $rd$ candidate representations.

Repeated exact certification is itself expensive, so it should be cached and
budgeted rather than skipped. Numerical approximations are useful as a
\emph{rejection prefilter}: clearly unequal candidates can often be discarded
cheaply. They should not be the final acceptance certificate.

\subsection{F12: absence of an application-level budget}
The outer loop tests its index against a stack whose length may grow. No
maximum number of trials, cumulative time limit, memory budget, polynomial
degree limit, or total candidate budget is enforced by the program. This does
not prove nontermination on every input; it means the caller lacks a
predictable resource contract.

Recommended controls include total trials and elapsed time, per-trial limits,
maximum candidate-generation size, a cap on allowed polynomial degrees, and
an explicit representation-size budget. Exhaustion should return a structured
``search incomplete'' result together with the best \emph{certified} value
found so far. It must not return an unverified intermediate as though it were
an exact simplification.

\section{Repair strategy and a conservative reference component}
\label{sec:repair}
\subsection{The order of repairs matters}
The first repair should preserve the original target, eliminate or guard
unsafe preprocessing, and make exact equality an acceptance requirement.
Improving speed before these changes merely makes an incorrect result easier
to obtain. The next repairs should reject invalid multipliers, guard solver
results and empty candidate sets, and restore session hygiene. Only then
should the multiplier heuristics be tuned and benchmarked.

For a wrapper over exact algebraic numbers, retain an immutable copy of the
original input before any factoring. Each local transformation should preserve
that value, and the final output should be certified against it. An exact
post-check can prevent a wrong result from escaping, but it does not make
unsafe transformations into a good algorithm: a failed certificate should
trigger fallback or a structured failure, not a second attempt to justify the
altered target.

\subsection{Separate value certification from representation quality}
Define a target expression language first. A radical-only language might
allow rational numbers, algebraic operations, and rational powers, but forbid
\code{Root} and \code{AlgebraicNumber} as output constructors. Alternatively,
an algebraic-canonicalization mode can allow those constructors explicitly.
These are different products and should have different success labels.

An example representation score, to be minimized lexicographically, is
\[
 \begin{aligned}
 \mathrm{score}(E)=(&\text{forbidden-constructor count},\ \raddepth(E),\\
                   &\text{radical-node count},\ \text{leaf count},\
                     \text{coefficient bit cost}).
 \end{aligned}
\]
This is a proposed engineering policy, not a theorem that one particular
score captures all mathematical simplicity. A rewrite is accepted only if it
is certified equal to the original and improves the selected score. Merely
returning a different expression should not count as improvement.

\subsection{A verified single-multiplier trial}
Appendix~\ref{app:trial} gives \code{VerifiedMultiplierTrial}, also supplied as
\code{verified_multiplier_trial.wl}. It deliberately accepts the root degree
$r$ as an explicit argument rather than inheriting the source's fragile depth
extraction. Its responsibilities are limited:
\begin{enumerate}
\item require an exact algebraic target and a nonzero exact algebraic
multiplier, with $r\ge2$;
\item compute the same minimal-polynomial/binomial GCD experiment;
\item enumerate phases and compare candidates to the \emph{original} target
using exact algebraic reduction;
\item return a named outcome or a controlled failure, with a per-call time
constraint.
\end{enumerate}
It does not claim completeness, minimal radical depth, a full memory budget,
or production validation. A result labeled \code{"VerifiedCandidate"} means
only that the returned value passed the exact equality check; an explicit
flag reports whether it contains \code{Root} or \code{AlgebraicNumber}.
The implementation sorts candidate representations before semantic
deduplication, so a first, unnecessarily large representative does not defeat
the intended leaf-count choice.

The code has been reviewed statically but was not executed in a Wolfram
kernel. It is a reference component illustrating the invariants, not a
silently substituted and supposedly tested replacement for all 585 lines of
the supplied program.

\subsection{Queue design and caching}
A revised search should keep the candidate generator separate from the trial
engine. Use a queue of named records, cache attempts by a certified canonical
multiplier key, and check the work budget before generating or enqueuing a
batch. Caching by \code{RootReduce[m]} is one possible policy for supported
exact algebraic numbers; the cost of canonicalization itself needs a budget.

Retain unresolved discriminant cofactors as metadata rather than discarding
them without a trace. A user can then distinguish ``all detected prime
combinations exhausted'' from ``factorization incomplete.'' The two statuses
have materially different meanings. Changes to search ordering should be
benchmarked independently of correctness, with the exact certification layer
left unchanged.

\section{Reproducibility and regression testing}
\label{sec:tests}
\subsection{Checks actually performed}
The archive contains \code{verify_algebra.py} and its generated
\code{test_results.json}. The script completed \textbf{45 independent checks}.
They include exact polynomial identities and GCDs for the examples, the
negative and complex branch calculations, counterexamples to the unsafe
power rules, five denominator inversions with exact remainders, the
polynomial-discriminant example, a degree-product example, the meaning of the
bad comparator, and exact candidate-count calculations.

One of the 45 checks is explicitly numerical rather than a proof: at 90-digit
working precision, 21 combinations of root degree and multiplier verified
that the phase enumeration contains a numerical representative of the
original target. The proof of that property is Theorem~\ref{thm:phase}, not the
numerical test. The JSON report distinguishes exact algebra, exact counting,
a source-semantics model of the comparator, and this numerical sanity check.

These are \emph{not 45 passing tests of the original program}. The Python
script neither executes nor claims to emulate Wolfram's evaluator, pattern
matcher, factorizer, or solver. No runtime benchmark, exact public-wrapper
transcript, kernel-specific message sequence, or original-code coverage
percentage is claimed.

\subsection{Tests prepared for a Wolfram kernel}
Appendix~\ref{app:tests} contains \code{run_wolfram_audit.wl}, which loads the
unmodified attachment in a disposable kernel and exercises 14 regression
contracts. Some are expected to fail on the original source: their purpose
is to expose defects. They check loading side effects, the comparator,
correct rationalization, direct and multiplier-assisted successes, original
branch preservation, unsafe factorization followed by substitution, zero
multiplier handling, an empty seed list, and a basic \code{alllevels} call.

The runner records the kernel version and restores the original
\code{Print} attributes on normal completion or an ordinary abort. It is not
a sandbox against a forcibly terminated process or all possible package side
effects. A fresh kernel remains important because the original file clears
several global definitions. The proposed reference component should be tested
separately; it is not loaded by this original-source regression runner.

\subsection{Further tests needed for a release}
A release-quality suite should add property-based tests generated from known
algebraic targets. Choose a simple exact $\beta$, construct $R=\beta^r$, keep
all branch metadata, and require any reported answer to equal $\beta$ exactly.
Vary signs and complex arguments, especially near principal-sector and branch
cut boundaries. Equality of $r$th powers is \emph{not} a sufficient oracle.

Test semantically equal but structurally different multipliers, rational and
Gaussian-rational denominators, mixed root degrees, near-zero nonzero algebraic
numbers, roots represented by \code{Root}, malformed options, approximate
inputs, and exhausted budgets. Test queue interruption and resumption with a
mock trial engine so that scheduling invariants can be checked without an
expensive field computation. Finally, test idempotence only as a chosen
simplification contract: failure to be idempotent can indicate an incomplete
normalization policy, but does not by itself prove unequal values.

\section{What should not be misdiagnosed}
Several unusual details are valid or at least not the defects they might
initially appear to be. One-argument \code{MinimalPolynomial} really does
return a callable pure-function representation. The minus sign and possible
nonmonic polynomial in the rationalization helper are correct. The
roots-of-unity enumeration correctly handles the multiplier phase when the
original target is preserved. A numerical sign decision followed by an exact
compensating sign does not alone change a product. Boolean non-strict sorting
predicates are supported; the actual prime-sort bug is the mistaken slot.

Likewise, a symbolic result from \code{Element} fails a \code{PatternTest}
rather than automatically being accepted. Several structurally different
candidates equal to one target are not several different algebraic values.
\code{Block} is not inherently an error, although it does not provide fresh
symbols or restore a separately changed protection attribute. A heuristic
returning the unchanged input is not a wrong-value result; it becomes
misleading only when presented as proof that no denesting exists.

These distinctions matter because repair effort should target demonstrated
failures and missing contracts rather than rewrite every unfamiliar idiom.

\section{Conclusion}
The program contains a substantial, mathematically motivated experiment:
multiplier selection can make a scaled root satisfy a smaller equation over
an algebraic coefficient field, and the polynomial-GCD step can expose that
fact. Its denominator-rationalization formula is also sound under explicit
nonzero algebraicity assumptions.

Its present correctness boundary is inadequate. The original branch can be
lost before candidate selection; the custom factorizer uses invalid general
power identities; the final zero test is heuristic; and several invalid
states are treated as though normal algebraic results were present. The
comparator and multiplier-count calculation introduce separate, concrete
search defects. Global session changes and implicit state layouts further
limit its suitability as reusable software.

The appropriate path forward is therefore \emph{certified heuristic
simplification}: preserve the original value, certify every accepted
candidate, measure representation improvement separately, and report search
limits honestly. Those changes retain the interesting mathematics while
removing the assumption that a plausible-looking radical expression must be
the same algebraic number.

\clearpage
\appendix
\section{Complete supplied source}
\label{app:source}
The following listing is the supplied attachment, with line numbers added by
LaTeX. The separate \code{original.wl} file preserves the original bytes.
\lstset{basicstyle=\ttfamily\fontsize{8}{9.5}\selectfont,numbers=left,
 numberstyle=\tiny\color{muted},numbersep=8pt,stepnumber=1,
 xleftmargin=20pt,backgroundcolor=\color{white},frame=none}
\begin{lstlisting}
ClearAll[Strad];
Unprotect[Print];
Strad[eaf__] := Block[{Print = (Null &)}, DenestRadicals3[eaf]];
ClearAll[DenestRadicals3];
DenestRadicals3[expr_, alllevels_ : False, func_ : denest11] :=
 Block[{f, mod, level, problems, i1, i2},
  If[alllevels =!= True,
   level = ReplaceAll,
   level = ReplaceRepeated
   ];
  mod = level[Factorc[expr],
    radicand_?AlgebraicQ^pow_Rational :> Block[{factored, SumQ, i, j},
      SumQ[expression_] :=
       MemberQ[
        expression, _Plus | Complex[Except[0], _], {0, Infinity}];
      factored = FactorTermsList[radicand];
      factored[[2]] = 
       Flatten[{Replace[Factorc[factored[[2]]], 
          prod_Times :> List @@ prod]}];
      factored[[1]] = 
       DeleteCases[{#, ComplexExpand[factored[[1]]/#]} &[
         GCD[Re[factored[[1]]], Im[factored[[1]]]]], 1];
      factored[[1]] = #[[1]]^#[[2]] & /@ 
        Flatten[FactorInteger /@ factored[[1]], 1];
      factored = Flatten[factored];
      factored = 
       factored //. 
        list : {a___, expr1_, b___, expr2_, c___} :> {a, b, c, 
           ComplexExpand[expr1*expr2]} /; And[SumQ[expr1], SumQ[expr2]];
      j = 1;
      For[{i = 1}, i <= Length[factored], i++,
       If[N[Re[factored[[i]]], 100] < 0,
        factored[[i]] = -factored[[i]];
        j *= -1
        ]
       ];
      factored = DeleteCases[Prepend[factored, j], 1];
      f[Simplify[
        Times @@ (Simplify[RationalizeDenominator[Together[#]]] & /@ 
           factored)], pow]
      ]
    ];
  mod = mod //. 
    prod : Times[a___, expr1 : f[rad1_, pow1_], b___, 
       expr2 : f[rad2_, pow2_], c___] :> 
     Times[a, b, c, f[(rad1^pow1)*(rad2^pow2)]] /; 
      And[! FreeQ[rad1, Power], ! FreeQ[rad2, Power]];
  mod = level[mod, f[rad_, pow_] :> f[rad^pow]];
  mod = mod //. 
    prod : Times[a___, expr1 : f[rad1_], b___, expr2 : f[rad2_], 
       c___] :> 
     Times[a, b, c, f[rad1*rad2]] /; 
      And[And[Re[#] != 0, Im[#] != 0] &[N[rad1, 10]], 
       And[Re[#] != 0, Im[#] != 0] &[N[rad2, 10]]];
  problems = DeleteDuplicates[Cases[mod, f[_], {0, Infinity}]];
  Print[problems];
  Print[];
  If[alllevels =!= True,
   mod = 
    Simplify[
     mod /. ((# -> Replace[#, f[val_] :> func[val]]) & /@ problems)],
   For[{i1 = 1}, MemberQ[problems, f[__]] === True, i1++,
    i1 = Mod[i1, Length[problems], 1];
    If[MatchQ[problems[[i1]], f[problem__] /; FreeQ[problem, f[__]]],
     problems[[i1]] = problems[[i1]] -> (problems[[i1]] /. f -> func);
     For[{i2 = 1}, i2 <= Length[problems], i2++,
      If[Head[problems[[i2]]] =!= Rule,
       problems[[i2]] = problems[[i2]] //. problems[[i1]];
       ]
      ];
     ];
    ];
   mod = Simplify[mod //. problems];
   ];
  Return[mod];
  ]
  
ClearAll[denest11];
denest11[problem_, forcedmultipliers_ : 0] :=
 Block[{radicand, outerpower, reductionpower, extrapower, 
   multiplierstack, multipliers, multiplier, terms, history, attempt, 
   solution, updatehistory, stackchange, addmultipliers, 
   trymultiplier, newmultipliers, x, i1, i2, i3},
  If[! MemberQ[problem, 
     Power[inner_?(Function[var, MemberQ[var, Power[a_, b_], Infinity]]),
       outer_], {0, Infinity}],
   Return[problem];
   ];
  Print["problem  ", problem];
  reductionpower = (LCM @@ (maxnestdepth[problem, 
        MatchQ[#, Power[_, _Rational]] &] /. {Power[base_, exp_], _} :>
         Denominator[exp]));
  outerpower = 1/reductionpower;
  radicand = problem^reductionpower;
  extrapower = 
   Times @@ ((LCM @@ #) & /@ (DeleteCases[
         Cases[nestdepthsort[radicand, 
           MatchQ[#, Power[_, _Rational]] &], 
          patt : {{expr_, nest_}, Repeated[{_, _}, {0, Infinity}]} /; 
           nest >= 2], {Except[Power[_, _Rational]], _}, {2}] /. 
        patt : {Power[base_, exp_], nest_} :> Denominator[exp]));
  Print["outer root: ", reductionpower, "  extra root: ", extrapower];
  Print[];
  
  (*tries a multiplier
  returns {<boolean progress>, {multiplier, minpoly, degree}, [answer]}*)
  trymultiplier[radicand_, outerpower_, multiplier_, best_, 
     worseknown_] /; MatchQ[outerpower, Rational[1, _Integer]] :=
   Block[{minpoly, degree, polygcd, possibilities, result, expected, 
     arg1},
    minpoly = MinimalPolynomial[(multiplier*radicand)^outerpower];
    degree = Exponent[minpoly[arg1], arg1];
    result = {False, {multiplier, minpoly, degree}};
    Print["multiplier: ", multiplier, "   minpoly degree: ", degree];
    If[Or[degree < best, GCD[degree, best] < best],
     polygcd = 
      PolynomialGCD[minpoly[arg1], 
       arg1^(1/outerpower) - multiplier*radicand, 
       Extension -> Automatic];
     If[And[Exponent[polygcd, arg1] < (1/outerpower), polygcd =!= 1],
      possibilities = 
       Flatten[Outer[Times, 
         arg1/(multiplier^outerpower) /. Solve[polygcd == 0, arg1], 
         Exp[2*I*\[Pi]*Range[0, (1/outerpower) - 1, 1]/(1/outerpower)]]];
      result = 
       DeleteDuplicates[
        Select[possibilities, PossibleZeroQ[# - radicand^outerpower] &]];
      If[Length[result] > 1,
       Print[Length[result], " possible solutions, taking the first"];
       Print[result];
       result = Take[result, 1];
       ];
      (*progress because a denesting occurred*)
      result = {True, {multiplier, minpoly, degree}, result[[1]]},
      If[polygcd =!= 1,
       If[
        Or[degree < best, 
         And[GCD[degree, best] < best, worseknown =!= True]],
        (*progress because degree of minpoly has decreased*)
        result = {True, {multiplier, minpoly, degree}}
        ],
       Print[
        "PolynomialGCD error: ", {radicand, outerpower, multiplier, 
         minpoly}]
       ]
      ],
     If[And[worseknown =!= True, degree != best],
      (*progress because this multiplier has greater degree than a pre\
vious multiplier*)
      Print["previous progress identified"];
      result = {True, {multiplier, minpoly, degree}};
      ]
     ];
    Return[result];
    ];
  
  (*form for attempts is {multiplier, minpoly, degree}, 
  attempt2 is Null if only one attempt is available*)
  newmultipliers[smaller_, larger_, reductionpower_] :=
   Block[{done, new, primes, multiplierprimes, minprimes, maxprimes, 
     primeratio, multipliercap, arg1},
    minprimes = 5;
    maxprimes = 14;
    primeratio = 100;
    multipliercap = 1000;
    
    (*two possibilites: reduction in degree, 
    reduction only in GCD of degrees*)
    If[Or[larger === Null, 
      GCD[smaller[[-1]], larger[[-1]]] == 
       Min[smaller[[-1]], larger[[-1]]]],
     (*one degree is a multiple of the other, usual case*)
     primes = 
      DeleteCases[
       FactorInteger[Discriminant[smaller[[2]][arg1], arg1], 
        Automatic], {-1, 1}];
     primes = 
      Sort[DeleteCases[primes, 
        factor : {prime_, _} /; 
         And[! PrimeQ[prime], (Print["composite factor: ", factor]; 
           True)]], #1[[1]] <= #[[2]] &];
     primes = Flatten[primes[[All, 1]]];
     multiplierprimes = primes;
     For[{i3 = Length[primes]}, 
      And[i3 > minprimes, 
       multiplierprimes[[i3]]/multiplierprimes[[i3 - 1]] > primeratio],
       i3--,
      multiplierprimes[[i3]] = 0;
      ];
     multiplierprimes = 
      DeleteCases[multiplierprimes, 0][[
       1 ;; Min[
         Max[minprimes, Ceiling[Log[multipliercap]/Log[reductionpower]]],
          Length[multiplierprimes]]]];
     new = 
      smaller[[1]]*
       Drop[Union[
         Divisors[(Times @@ multiplierprimes)^(reductionpower - 1)], 
         primes], 1];
     ,
     (*different degree reductions in each degree from a larger degree\
*)
     Print[
      "GCD of smallest degrees is smaller than any observed degree: \
", {smaller[[{1, -1}]], larger[[{1, -1}]]}];
     new = smaller[[1]]*larger[[1]]*{1};
     Print[PolynomialGCD @@ {smaller[[2]][x], larger[[2]][x]}];
     ];
    new = {{False, smaller}, ComplexitySort[new]};
    Return[new];
    ];
  
  (*compute initial set of multipliers*)
  If[ListQ[forcedmultipliers],
   multipliers = {{False, {0, Infinity, Infinity}}, forcedmultipliers},
   (*old, may be better*)
   (*
   terms=DeleteCases[Replace[#,{Times[a___,
   Alternatives[_Rational,_Integer],b___]\[RuleDelayed]a*b,
   Alternatives[_Rational,_Integer]\[Rule]0,Complex[a_,
   b_]\[RuleDelayed]Sign[b]*I}]&/@(List@@Expand[
   RationalizeDenominator[radicand]]),0];
   *)
   terms = 
    DeleteCases[
     Replace[#, {Times[a___, Alternatives[_Rational, _Integer], 
           b___] :> a*b, Alternatives[_Rational, _Integer] -> 0, 
         Complex[a_, b_] :> Sign[b]*I}] & /@ (List @@ 
        Expand[RationalizeDenominator[radicand]]), 0];
   multipliers = Prepend[Replace[If[Head[#] === Times,
         
         Times @@ (Function[{var}, 
             var^(maxnestdepth[var, 
                  Function[{var2}, 
                   MatchQ[var2, Power[_, _Rational]]]][[
                 1]] /. {{Power[base_, 
                    exp_], _} :> (Denominator[exp] - Numerator[exp])/
                   Numerator[exp], {Complex[
                    0, _], _} :> -1})] /@ (List @@ #)),
         #^(maxnestdepth[#, 
              Function[{var}, MatchQ[var, Power[_, _Rational]]]][[
             1]] /. {{Power[base_, 
                exp_], _} :> (Denominator[exp] - Numerator[exp])/
               Numerator[exp], {Complex[0, _], _} :> -1})
         ], 
        Times[a___, Alternatives[_Rational, _Integer], b___] :> 
         a*b] & /@ Sort[terms, (Re[#1] <= Re[#2]) &], 1];
   multipliers = {{True, {0, Infinity, Infinity}}, multipliers};
   ];
  
  solution = problem;
  history = {{0, Infinity, Infinity}};
  (*multiplierstack form {,{<multipliers>}},...}*)
  multiplierstack = {multipliers};
  (*add power set of multipliers as second element in multiplierstack*)
  (*multiplierstack=Append[multiplierstack,
  {multiplierstack[[1,1]],ComplexitySort[DeleteCases[Union[(Times@@#)&/@
  Subsets[DeleteCases[multipliers[[2]],1]]],mult_/;MemberQ[
  multipliers[[2]],mult]]]}];*)
  Catch[
   For[{i1 = 1}, i1 <= Length[multiplierstack], i1++,
    (*compute new multipliers from data if not already done*)
    If[ListQ[multiplierstack[[i1, 2]]] =!= True,
     multiplierstack[[i1]] = 
       multiplierstack[[i1, -2]] @@ (multiplierstack[[i1, -1]]);
     ];
    multipliers = multiplierstack[[i1]];
    stackchange = False;
    If[multipliers[[1, 2]] =!= history[[1]],
     Print["multiplier count: ", Length[multipliers[[2]]], "  from: ",
       multipliers[[1, 2, {1, -1}]]],
     If[! ListQ[forcedmultipliers],
      Print["multiplier count: ", Length[multipliers[[2]]], 
        "  multipliers: ", multipliers[[2]]];,
      Print["multiplier count: ", Length[multipliers[[2]]]];
      ]
     ];
    For[{i2 = 1}, i2 <= Length[multipliers[[2]]], i2++,
     multiplier = multipliers[[2, i2]];
     updatehistory = False;
     addmultipliers = False;
     (*try a multiplier, different based on history*)
     If[And[Length[history] >= 3, history[[-2]] =!= history[[1]]],
      attempt = 
       trymultiplier[radicand, outerpower, multiplier, 
        history[[-1, -1]], True],
      attempt = 
        trymultiplier[radicand, outerpower, multiplier, 
         history[[-1, -1]], False];
      ];
     multiplierstack[[i1, 2, i2]] = {multiplier, attempt[[2, -1]]};
     
     If[And[multipliers[[1, 1]] === True, 
       attempt[[2, -1]] <= history[[-1, -1]]],
      addmultipliers = True
      ];
     (*check if progress has been made on this attempt, 
     no action taken if no progress*)
     If[attempt[[1]] === True,
      (*history will be updated*)
      updatehistory = True;
      
      (*decide if a stackchange will be made*)
      If[history[[-1]] =!= history[[1]],
       stackchange = True;
       ];
      
      If[ListQ[attempt[[-1]]] =!= True,
       (*a solution was obtained*)
       Print["success: ", 
        Join[attempt[[2, {1, -1}]], 
         If[IntegerQ[multiplier], {FactorInteger[multiplier]}, {}]]];
       solution = Simplify[attempt[[-1]]];
       If[
        Or[history[[-1]] === history[[1]], 
         history[[2, -1]]/attempt[[2, -1]] >= reductionpower],
        (*not worth trying for additional denesting, exit immediately*)
        history = Append[history, attempt[[2]]];
        Throw[, "solution"],
        (*further denesting will be attempted*)
        Print["attempting additional denesting, current solution is:"];
        Print[solution];
        ];
       ];
      
      If[
       And[history[[-1]] =!= history[[1]], 
        Or[attempt[[2, -1]] < history[[-1, -1]], 
         history[[-2]] === history[[1]]]],
       (*improvement has occured, new multipliers will be created*)
       Print["improvement  ", attempt[[2, {1, -1}]]];
       addmultipliers = True;
       ];
      ];
     
     (*add new multipliers if needed*)
     If[addmultipliers === True,
      Block[{smaller, larger, done, new},
       If[attempt[[2, -1]] <= history[[-1, -1]],
        smaller = attempt[[2]];
        larger = history[[-1]],
        smaller = history[[-1]];
        larger = attempt[[2]];
        ];
       new = {{False, smaller}, 
         newmultipliers, {smaller, 
          If[larger =!= history[[1]], larger, Null], reductionpower}};
       
       Print["adding new multipliers from: ", multiplier];
       (*update multiplierstack with new multipliers*)
       If[multipliers[[1, 1]] === False,
        (*use new multipliers immediately*)
        done = {multipliers[[1]], multipliers[[2, 1 ;; i2]]};
        multiplierstack = Insert[multiplierstack, new, i1 + 1];
        If[i2 < Length[multipliers[[2]]],
         
         multiplierstack = 
           Insert[multiplierstack, {multipliers[[1]], 
             multipliers[[2, i2 + 1 ;; -1]]}, i1 + 2];
         ];
        multiplierstack[[i1]] = done,
        (*use new multipliers later*)
        stackchange = False;
        For[{i3 = Length[multiplierstack]}, i3 >= i1, i3--,
         
         If[Or[multiplierstack[[i3, 1, 2]] === history[[1]], 
           new[[1, 2, -1]] > multiplierstack[[i3, 1, 2, -1]], 
           And[new[[1, 2, -1]] === multiplierstack[[i3, 1, 2, -1]], 
            LeafCount[new[[1, 2, 1]]] >= 
             LeafCount[multiplierstack[[i3, 1, 2, 1]]]]],
          multiplierstack = Insert[multiplierstack, new, i3 + 1];
          Break[];
          ]
         ];
        ];
       ]
      ];
     
     (*update history if needed*)
     If[updatehistory === True,
      For[{i3 = Length[history]}, i3 >= 1, i3--,
        If[history[[i3, -1]] > attempt[[2, -1]],
         history = Insert[history, attempt[[2]], i3 + 1];
         Break[];
         ]
        ];
      ];
     
     (*break inner loop if stack has changed*)
     If[stackchange === True,
      Break[];
      ];
     ];(*end of inner loop*)
    
    (*add additional multipliers after first pass (if no forced multip\
liers)*)
    If[And[i1 === 1, ! ListQ[forcedmultipliers]],
     Block[{best, worst, extra, f},
      best = Min[multiplierstack[[i1, 2, All, 2]]];
      worst = Max[multiplierstack[[i1, 2, All, 2]]];
      extra = 
       DeleteCases[
        DeleteDuplicates[
         Replace[(Distribute[
               f @@ (#^Range[0, 1] & /@ (Cases[
                    multiplierstack[[i1, 2]], {mult_, degree_} /; 
                    degree < worst][[All, 1]])), 
               List] //. {f[a___, Power[base_Times, exp_], b___] :> 
                f[a, b, Sequence @@ ((List @@ base)^exp)], 
               f[a___, prod_Times, b___] :> f[a, b, Sequence @@ prod],
                f[a___, Power[base_, exp_], b___] :> 
                f[a, b, (-1)^exp, (-base)^exp] /; 
                 And[base != -1, Re[base] < 0], 
               f[a___, Power[base_, Rational[num1_, denom_]], b___, 
                 Power[base_, Rational[num2_, denom_]], c___] :> 
                f[a, b, c, base^(Mod[num1 + num2, denom]/denom)]}) /. 
            f -> Times, Times[a___, _Integer, b___] :> a*b, {1}] //. 
          Power[base_, Rational[a_, b_]] :> base^(Mod[a, b]/b)], 
        Alternatives @@ multiplierstack[[i1, 2, All, 1]]];
      If[Length[extra] >= 1,
       Print["adding multipliers based on progress in initial guess"];
       multiplierstack = 
        Insert[multiplierstack, {multiplierstack[[i1, 1]], extra}, 
         i1 + 1];
       ];
      ]
     ];
    
    ](*end of outer loop*),
   "solution"
   ];(*end of Catch, used to break the loop when a solution is found*)
  
  Print["history: ", #[[{1, -1}]] & /@ history];
  If[solution =!= problem,
   Print[solution],
   Print["no progress"];
   ];
  Print[];
  Return[solution];
  ]
  
ClearAll[RationalizeDenominator];
RationalizeDenominator[expr_] :=
 Block[{conjugate, denom, x},
  denom = Denominator[expr];
  conjugate = MinimalPolynomial[denom];
  Expand[
    Numerator[expr]*PolynomialQuotient[conjugate[x], x - denom, x] /. 
     x -> 0]/(-1*conjugate[0])
  ]
  
ClearAll[AlgebraicQ];
AlgebraicQ[num_] := Element[num, Algebraics]

ClearAll[nestdepths]
nestdepths[expr_, func_] :=
 Block[{inner},
  If[Depth[expr] === 1,
   {expr, If[func[expr], 1, 0]},
   inner = Map[nestdepths[#, func] &, (List @@ expr), {1}];
   {expr, 
    Sequence @@ 
     inner, ((Max @@ ((#[[-1]]) & /@ inner)) + If[func[expr], 1, 0])}
   ]
  ]

ClearAll[nestdepthsort]
nestdepthsort[expr_, func_] :=
 Block[{data},
  data = nestdepths[expr, func];
  data = 
   data //. 
    patt : {a : Except[_List], b : __, c : Except[_List]} :> {{a, c}, 
      b};
  data = data //. patt : {rep : Repeated[{_, _}]} :> Sequence @@ patt;
  data = Gather[{data}, (#1[[2]] == #2[[2]]) &];
  data = Sort[data, (#1[[1, 2]] >= #2[[1, 2]]) &];
  data
  ]

ClearAll[maxnestdepth]
maxnestdepth[expr_, func_] :=
 Block[{data, max},
  data = nestdepths[{expr}, func][[2]];
  max = data[[-1]];
  data = 
   FixedPoint[
    DeleteCases[# /. {a_, b_} :> {} /; b < max, {}, Infinity] &, data];
  data = 
   If[And[ListQ[#], Length[#] > 2], Sequence @@ (#[[2 ;; -2]]), #] & /@
     data;
  If[Length[data] > 2,
   data = data[[2 ;; -2]],
   data = {data}
   ];
  data
  ]
  
ClearAll[ComplexitySort];
ComplexitySort[expr_] := 
  Return[Sort[expr, 
    Switch[LeafCount[#1] - LeafCount[#2], a_ /; a < 0, True, 
      a_ /; a > 0, False, 0, Abs[#1] <= Abs[#2]] &]];
	  
ClearAll[Factorc]
Factorc[expr_] :=
 Flatten[
   fullfactor[
     expr] //. {pow : {a : Except[_List], b : Except[_List]} :> {a^b},
      prod : {rep : (Repeated[{Except[_List]}, {2, 
            Infinity}])} :> {Times[rep]}, 
     sum : {rep : Repeated[{{Except[_List]}}, {2, Infinity}]} :> 
      Plus[rep], 
     pow2 : {{a : Except[_List]}, b : Except[_List]} :> {a^b}, 
     pow3 : {{{a : Except[_List]}}, b : Except[_List]} :> {a^b}}][[1]]

ClearAll[fullfactor]
fullfactor[expr_] :=
 Block[{factorlist, i = 0, j = 0, done = False},
  factorlist = FactorList[expr];
  While[! done && j < 5,
   j++;
   done = True;
   For[{i = 1}, i <= Length[factorlist], i++,
    Which[
     Head[factorlist[[i, 1]]] === Power,
     done = False;
     factorlist[[
       i]] = {{factorlist[[i, 1, 1]], 
        factorlist[[i, 1, 2]]*factorlist[[i, 2]]}},
     And[
      MemberQ[{Integer, Rational}, Head[factorlist[[i, 1]]]], ! 
       PrimeQ[factorlist[[i, 1]]], Abs[factorlist[[i, 1]]] =!= 1],
     done = False;
     factorlist[[i]] = {#[[1]], #[[2]]*factorlist[[i, 2]]} & /@ 
       FactorInteger[factorlist[[i, 1]]],
     Head[factorlist[[i, 1]]] === Plus,
     done = False;
     factorlist[[i]] = 
      combine[fullfactor /@ 
        Replace[factorlist[[i, 1]], Plus[a_, b__] :> {a, b}], 
       factorlist[[i, 2]]],
     True,
     factorlist[[i]] = {factorlist[[i]]}
     ]
    ];
   factorlist = Flatten[factorlist, 1]
   ];
  factorlist
  ]

ClearAll[combine];
combine[terms_List, pow_] :=
 Block[{simpler = terms, base, bases, positions, factor, i, j, k},
  For[{i = 1}, i <= Length[terms], i++,
   base = terms[[i]];
   For[{j = 1}, j < Length[base], j++,
    For[{k = j + 1}, k <= Length[base], k++,
     If[base[[j, 1]] === base[[k, 1]],
      base[[j]] = {base[[j, 1]], base[[j, 2]] + base[[k, 2]]};
      base[[k]] = {0, 0}
      ]
     ];
    If[base[[j, 1]] === 1,
     base[[j]] = {0, 0}
     ]
    ];
   simpler[[i]] = Delete[base, Position[base, {0, 0}]];
   ];
  bases = Intersection @@ Map[#[[1]] &, simpler, {2}];
  For[{i = 1}, i <= Length[bases], i++,
   positions = 
    Position[#, {bases[[i]], Except[_List]}, 1] & /@ simpler;
   factor = 
    Min[Extract[simpler[[#]], positions[[#]]][[1, 2]] & /@ 
      Range[1, Length[simpler], 1]];
   For[{j = 1}, j <= Length[simpler], j++,
    simpler[[j, Sequence @@ (positions[[j, 1]]), 2]] -= factor
    ];
   bases[[i]] = {bases[[i]], factor*pow};
   ];
  Append[bases, {simpler, pow}]
  ]
\end{lstlisting}

\clearpage
\section{Defensive single-trial reference code}
\label{app:trial}
This reference component is not a complete denester and was not executed in a
Wolfram kernel during the review. It deliberately calls a certified value a
``candidate,'' not a certified reduction of radical depth.
\lstset{basicstyle=\ttfamily\fontsize{8.5}{10}\selectfont,firstnumber=1}
\begin{lstlisting}
(* A conservative, single-multiplier reference component for the article.
   This is NOT a complete replacement denester, and was not executed in a
   Wolfram kernel during this review. The original target is never discarded.
   Load this file separately; it does not load or modify the supplied code. *)
BeginPackage["RadicalAudit`"];
VerifiedMultiplierTrial::usage =
  "VerifiedMultiplierTrial[beta, r, m] tries one nonzero exact algebraic multiplier, preserving the original branch. A verified candidate is not a certificate of radical-depth reduction.";
Begin["`Private`"];

ClearAll[exactAlgebraicValueQ, exactEqualQ];
exactAlgebraicValueQ[z_] :=
  TrueQ[NumericQ[z] && Precision[z] === Infinity &&
    Element[z, Algebraics]];
exactEqualQ[a_, b_] :=
  TrueQ[Quiet[Check[RootReduce[a - b] === 0, False]]];

Options[VerifiedMultiplierTrial] = {TimeConstraint -> 30};
VerifiedMultiplierTrial[beta_, r_, m_, OptionsPattern[]] :=
 Module[{limit = OptionValue[TimeConstraint]},
  If[!(limit === Infinity ||
      TrueQ[NumericQ[limit] && Im[limit] == 0 && limit > 0]),
   Return[Failure["InvalidTimeConstraint", <||>]]];
  TimeConstrained[
   Quiet[Check[
     Module[{x, radicand, scaled, minpoly, g, d, rules,
       roots, phases, candidates, valid, best},
      If[!TrueQ[IntegerQ[r] && r >= 2],
       Return[Failure["InvalidRootDegree", <|"Degree" -> r|>]]];
      If[!exactAlgebraicValueQ[beta] || !exactAlgebraicValueQ[m],
       Return[Failure["ExactAlgebraicInputsRequired", <||>]]];
      If[exactEqualQ[m, 0],
       Return[Failure["ZeroMultiplier", <||>]]];
      If[exactEqualQ[beta, 0],
       Return[<|"Status" -> "Unchanged", "Value" -> 0,
         "Verified" -> True|>]];
      radicand = beta^r;
      scaled = m radicand;
      minpoly = MinimalPolynomial[scaled^(1/r), x];
      If[!PolynomialQ[minpoly, x],
       Return[Failure["MinimalPolynomialFailed", <||>]]];
      g = PolynomialGCD[minpoly, x^r - scaled,
        Extension -> Automatic];
      d = Exponent[g, x];
      If[!PolynomialQ[g, x] || !IntegerQ[d],
       Return[Failure["PolynomialGCDFailed", <||>]]];
      If[!(1 <= d < r),
       Return[<|"Status" -> "NoProperGCD", "Value" -> beta,
         "GCDDegree" -> d|>]];
      rules = Solve[g == 0, x];
      If[!ListQ[rules] || rules === {} ||
        !TrueQ[And @@ (MatchQ[#, {_Rule}] & /@ rules)],
       Return[Failure["UnexpectedSolveResult", <||>]]];
      roots = x /. rules;
      phases = Table[Exp[2 Pi I k/r], {k, 0, r - 1}];
      candidates = Flatten[
        Outer[Times, roots/m^(1/r), phases], 1];
      valid = Select[candidates,
        exactAlgebraicValueQ[#] && exactEqualQ[#, beta] &];
      valid = SortBy[valid, LeafCount];
      valid = DeleteDuplicates[valid, exactEqualQ];
      If[valid === {},
       Return[Failure["NoCertifiedCandidate", <||>]]];
      best = First[valid];
      <|"Status" -> "VerifiedCandidate", "Original" -> beta,
        "Value" -> best, "Verified" -> True,
        "Multiplier" -> m, "GCDDegree" -> d,
        "ContainsRootObjects" ->
          !FreeQ[best, _Root | _AlgebraicNumber]|>
      ],
     Failure["EvaluationFailed", <||>]]],
   limit, Failure["TimeLimit", <||>]]
  ];
End[];
EndPackage[];
\end{lstlisting}

\clearpage
\section{Unexecuted Wolfram regression contracts}
\label{app:tests}
Run these tests on the original source in a fresh kernel. Some tests are
expected to fail; they express desired correctness and error-handling
contracts, not a claim that the original implementation satisfies them.
\lstset{firstnumber=1}
\begin{lstlisting}
(* Regression contracts for the UNMODIFIED supplied source.
   NOT EXECUTED in a Wolfram kernel during this review.
   Some tests are EXPECTED TO FAIL: they expose defects, rather than documenting
   correct existing behavior. Run only in a fresh/disposable kernel because
   original.wl clears global definitions and unprotects System`Print.
   Invocation: wolframscript -file run_wolfram_audit.wl
   The original Print attributes are restored on normal completion or Abort. *)
ClearAll[auditEqualQ, auditEval];
auditEqualQ[a_, b_] :=
  TrueQ[Quiet[Check[RootReduce[a - b] === 0, False]]];
SetAttributes[auditEval, HoldAll];
auditEval[expr_] :=
  TimeConstrained[Block[{$Output = {}}, Quiet[expr]], 20, $TimedOut];

Module[{dir, savedPrintAttributes, report, protectionOK},
 dir = DirectoryName[$InputFileName];
 savedPrintAttributes = Attributes[Print];
 CheckAbort[
  Get[FileNameJoin[{dir, "original.wl"}]];
  protectionOK = MemberQ[Attributes[Print], Protected];
  report = TestReport[{
    VerificationTest[protectionOK, True,
      TestID -> "loading-preserves-Print-protection"],
    VerificationTest[(#1[[1]] <= #[[2]] &)[{2, 1}, {3, 1}], True,
      TestID -> "prime-comparator-orders-2-before-3"],
    VerificationTest[
      auditEqualQ[RationalizeDenominator[1/(1 + Sqrt[2])],
        Sqrt[2] - 1], True, TestID -> "rationalizer-real"],
    VerificationTest[
      auditEqualQ[RationalizeDenominator[1/(1 + I)], (1 - I)/2],
      True, TestID -> "rationalizer-complex"],
    VerificationTest[
      auditEqualQ[auditEval[denest11[Sqrt[3 + 2 Sqrt[2]], {1}]],
        1 + Sqrt[2]], True, TestID -> "direct-denesting-success"],
    VerificationTest[
      auditEqualQ[auditEval[denest11[Sqrt[5 + 2 Sqrt[6]], {2}]],
        Sqrt[2] + Sqrt[3]], True,
      TestID -> "rational-multiplier-success"],
    VerificationTest[
      auditEqualQ[auditEval[denest11[-Sqrt[3 + 2 Sqrt[2]], {1}]],
        -1 - Sqrt[2]], True, TestID -> "negative-original-branch"],
    VerificationTest[
      auditEqualQ[auditEval[denest11[
        Sqrt[-1 + 2 I Sqrt[2]] Sqrt[-2 + 2 I Sqrt[3]], {1}]],
        (1 + I Sqrt[2]) (1 + I Sqrt[3])],
      True, TestID -> "grouped-complex-original-branch"],
    VerificationTest[
      auditEqualQ[auditEval[DenestRadicals3[
        Sqrt[-1 + 2 I Sqrt[2]] Sqrt[-2 + 2 I Sqrt[3]]]],
        (1 + I Sqrt[2]) (1 + I Sqrt[3])],
      True, TestID -> "public-wrapper-complex-product"],
    VerificationTest[
      Module[{t}, auditEqualQ[
        auditEval[Factorc[Sqrt[t^2]]] /. t -> -2, 2]],
      True, TestID -> "nested-power-symbolic-branch"],
    VerificationTest[
      Module[{t}, auditEqualQ[
        auditEval[Factorc[Sqrt[t + t^2]]] /. t -> -2, Sqrt[2]]],
      True, TestID -> "factored-sum-symbolic-branch"],
    VerificationTest[
      MatchQ[auditEval[denest11[Sqrt[3 + 2 Sqrt[2]], {0}]], _Failure],
      True, TestID -> "zero-multiplier-controlled-rejection"],
    VerificationTest[
      auditEqualQ[auditEval[denest11[Sqrt[3 + 2 Sqrt[2]], {}]],
        Sqrt[3 + 2 Sqrt[2]]], True, TestID -> "empty-seed-list-no-change"],
    VerificationTest[
      auditEqualQ[auditEval[DenestRadicals3[
        Sqrt[3 + 2 Sqrt[2]], True]], 1 + Sqrt[2]],
      True, TestID -> "alllevels-basic-correctness"]
    }];
  Attributes[Print] = savedPrintAttributes;
  Print[report];
  Print["Kernel version: ", $Version];,
  Attributes[Print] = savedPrintAttributes;
  Abort[]
  ]
 ];
\end{lstlisting}

\clearpage
\begin{thebibliography}{99}
\bibitem{source} Supplied attachment, \emph{Pasted text.txt}, reproduced as
\code{original.wl}; 585 physical lines, 21,199 bytes. Source-line citations in
this article refer to this file. No authorship, date of original composition,
or supported Wolfram version was supplied.
\bibitem{power} Wolfram Research, \emph{Power}, Wolfram Language documentation.
\url{https://reference.wolfram.com/language/ref/Power.html}.
\bibitem{minimalpolynomial} Wolfram Research, \emph{MinimalPolynomial}, Wolfram
Language documentation.
\url{https://reference.wolfram.com/language/ref/MinimalPolynomial.html}.
\bibitem{polynomialgcd} Wolfram Research, \emph{PolynomialGCD}, Wolfram Language
documentation.
\url{https://reference.wolfram.com/language/ref/PolynomialGCD.html}.
\bibitem{possiblezero} Wolfram Research, \emph{PossibleZeroQ}, Wolfram Language
documentation, including the exact-algebraics method and documented false
positive examples.
\url{https://reference.wolfram.com/language/ref/PossibleZeroQ.html}.
\bibitem{rootreduce} Wolfram Research, \emph{RootReduce}, Wolfram Language
documentation.
\url{https://reference.wolfram.com/language/ref/RootReduce.html}.
\bibitem{factorinteger} Wolfram Research, \emph{FactorInteger}, Wolfram Language
documentation, including the partial-factorization \code{Automatic} argument.
\url{https://reference.wolfram.com/language/ref/FactorInteger.html}.
\bibitem{sort} Wolfram Research, \emph{Sort}, Wolfram Language documentation.
\url{https://reference.wolfram.com/language/ref/Sort.html}.
\bibitem{function} Wolfram Research, \emph{Function}, Wolfram Language
documentation.
\url{https://reference.wolfram.com/language/ref/Function.html}.
\bibitem{block} Wolfram Research, \emph{Block}, Wolfram Language documentation.
\url{https://reference.wolfram.com/language/ref/Block.html}.
\bibitem{module} Wolfram Research, \emph{Module}, Wolfram Language documentation.
\url{https://reference.wolfram.com/language/ref/Module.html}.
\bibitem{unprotect} Wolfram Research, \emph{Unprotect}, Wolfram Language
documentation.
\url{https://reference.wolfram.com/language/ref/Unprotect.html}.
\bibitem{replacerepeated} Wolfram Research, \emph{ReplaceRepeated}, Wolfram
Language documentation.
\url{https://reference.wolfram.com/language/ref/ReplaceRepeated.html}.
\bibitem{factorlist} Wolfram Research, \emph{FactorList}, Wolfram Language
documentation.
\url{https://reference.wolfram.com/language/ref/FactorList.html}.
\bibitem{element} Wolfram Research, \emph{Element}, Wolfram Language
documentation.
\url{https://reference.wolfram.com/language/ref/Element.html}.
\bibitem{patterntest} Wolfram Research, \emph{PatternTest}, Wolfram Language
documentation.
\url{https://reference.wolfram.com/language/ref/PatternTest.html}.
\bibitem{complexexpand} Wolfram Research, \emph{ComplexExpand}, Wolfram Language
documentation.
\url{https://reference.wolfram.com/language/ref/ComplexExpand.html}.
\bibitem{solve} Wolfram Research, \emph{Solve}, Wolfram Language documentation,
including \code{Cubics} and default algebraic-root representations.
\url{https://reference.wolfram.com/language/ref/Solve.html}.
\end{thebibliography}
\noindent\small All external documentation above was consulted on September 5,
2026. It establishes language semantics; the source-specific conclusions,
proofs, and proposed changes are the analysis in this article.
\end{document}
